
\documentclass[12pt]{article}
\pdfoutput=1

% =====================
% Packages and settings
% =====================
\usepackage{amsmath,amssymb,amsthm,amsfonts,bm,physics,mathrsfs,upgreek,bbm}
\usepackage{graphicx}
\usepackage{float}
\usepackage{enumitem}
\usepackage{array}
\usepackage{booktabs}
\usepackage{multirow}
\usepackage{afterpage}
\usepackage{listings}
\usepackage{subcaption}
\usepackage{algorithm}
\usepackage{algorithmic}
\usepackage{changepage}

\newcommand{\INPUT}{\item[\textbf{Input:}]}

\usepackage[hang,flushmargin]{footmisc}
\usepackage{setspace}
\usepackage{xcolor}
\usepackage[hypertexnames=false]{hyperref}
\usepackage{xr-hyper}
\usepackage[capitalize,noabbrev]{cleveref}
\usepackage[round]{natbib}
\setlength{\bibsep}{0pt}
\definecolor{Bleu}{RGB}{0,0,204}
\hypersetup{
  colorlinks,
  citecolor=Bleu,
  linkcolor=Bleu,
  urlcolor=Bleu
}
\makeatletter
\newcommand*{\addFileDependency}[1]{%
  \typeout{(#1)}%
  \@addtofilelist{#1}%
  \IfFileExists{#1}{}{\typeout{No file #1.}}%
}
\makeatother

\newcommand*{\myexternaldocument}[1]{%
  \externaldocument{#1}%
  \addFileDependency{#1.tex}%
  \addFileDependency{#1.aux}%
}

\myexternaldocument{main_jasa_appendix}

% =====================
% Math and theorem defs
% =====================
\DeclareMathOperator*{\argmin}{argmin}
\DeclareMathOperator*{\argmax}{argmax}
\DeclareMathOperator*{\esssup}{ess\,sup}
\DeclareMathOperator*{\essinf}{ess\,inf}
\renewcommand{\P}{\mathbb{P}}

\newcommand{\blue}[1]{{\color{blue}#1}}
\newcommand{\addTK}{\widetilde T_{\mathcal K}}
\newcommand{\adk}{\widetilde{\partial}_k}
\newcommand\independent{\protect\mathpalette{\protect\independenT}{\perp}}
\def\independenT#1#2{\mathrel{\rlap{$#1#2$}\mkern2mu{#1#2}}}

\DeclareFontFamily{U}{jkpmia}{}
\DeclareFontShape{U}{jkpmia}{m}{it}{<->s*jkpmia}{}
\DeclareFontShape{U}{jkpmia}{bx}{it}{<->s*jkpbmia}{}
\DeclareMathAlphabet{\mathfrak}{U}{jkpmia}{m}{it}
\SetMathAlphabet{\mathfrak}{bold}{U}{jkpmia}{bx}{it}

\theoremstyle{plain}
\newtheorem{theorem}{Theorem}
\newtheorem{lemma}[theorem]{Lemma}
\newtheorem{corollary}{Corollary}
\newtheorem{definition}{Definition}
\newtheorem{assumption}{Assumption}
\newtheorem{condition}{Condition}
\newtheorem{example}{Example}

% Spacing + blind toggle
\def\spacingset#1{\renewcommand{\baselinestretch}{#1}\small\normalsize}
\newcommand{\blind}{0}

% JASA margin rules
\addtolength{\oddsidemargin}{-.5in}
\addtolength{\evensidemargin}{-1in}
\addtolength{\textwidth}{1in}
\addtolength{\textheight}{1.87in}
\addtolength{\topmargin}{-1in}

% Abstract keyword command
\providecommand{\keywords}[1]{\small{\textbf{\textit{Keywords---}} #1}}

% =====================
% Document begins
% =====================
\usepackage[compact]{titlesec}
 

\newcommand{\JASATitleText}{Efficient Inference for Inverse Reinforcement Learning and Dynamic Discrete Choice Models}
\title{\bf \JASATitleText}
\author{
  Lars van der Laan\thanks{Corresponding author: lvdlaan@uw.edu} \\
  University of Washington \\
  \and
  Aur\'elien Bibaut \\
  Netflix Research \\
  \and
  Nathan Kallus \\
  Cornell Tech, Cornell University
}
\spacingset{1}




 
\begin{document}

\if0\blind
{
  \maketitle
} \fi

\if1\blind
{
  \bigskip
  \bigskip
  \bigskip
  \begin{center}
    {\LARGE\bf \JASATitleText}
  \end{center}
  \medskip
} \fi

\bigskip
\begin{abstract}
Inverse reinforcement learning (IRL) and dynamic discrete choice (DDC) models
explain sequential decision-making through latent reward functions. Existing
flexible IRL methods allow rich reward representations but typically do not
support valid inference, whereas classical DDC approaches provide inference only
under restrictive parametric structure. We develop a semiparametric framework
for \emph{debiased inverse reinforcement learning} in maximum-entropy IRL and
Gumbel-shock DDC models. Our key identifying observation is that the
log-behaviour policy can be treated as a pseudo-reward. It point-identifies
policy value differences and, under a normalization constraint, the reward
itself. This lets us express reward-dependent estimands as smooth functionals of
the behaviour policy and transition kernel. We establish pathwise
differentiability,
derive efficient influence functions, and construct automatic debiased
machine-learning estimators that permit flexible nuisance estimation while
attaining \(\sqrt{n}\)-consistency, asymptotic normality, and semiparametric
efficiency. The result is a computationally tractable framework connecting
modern IRL with classical DDC inference.
\end{abstract}

\noindent
\keywords{inverse reinforcement learning, dynamic discrete choice, semiparametric inference, debiased machine learning, policy evaluation}
\newpage
\spacingset{1.9}  % JASA spacing rule


\section{Introduction}
\label{sec:intro}
Behavioural data arise in robotics, human--computer interaction, healthcare,
and economics. Inverse reinforcement learning (IRL) seeks to recover the
latent reward or utility function underlying observed sequential decisions. It
has been used to model expert driving and robotic control
\citep[e.g.][]{abbeel2004apprenticeship,wulfmeier2015deep},
user behaviour in interactive systems
\citep[e.g.][]{chandramohan2012user,hossain2023bayesian},
and clinical decision-making in critical care
\citep[e.g.][]{yu2019inverse,bovenzi2025identifying}. Because rewards summarize preferences or tradeoffs, they provide a basis for
counterfactual analysis---for example, predicting how behaviour would change
under new policies, constraints, or environments---that is not available from
policy prediction alone.

A central difficulty is that IRL is inherently ill posed: many rewards can
rationalize the same behavior
\citep{ng1999policy,ng2000algorithms,cao2021identifiability,
skalse2023invariance,skalse2024partial}. To make the problem tractable, IRL and
dynamic discrete choice (DDC) models impose an optimality model linking rewards
to behavior. Exact optimality implies deterministic policies and is often too
rigid for observed human behavior. A foundational stochastic alternative in DDC
adds Type-I extreme-value utility shocks, yielding choice probabilities of
softmax form and a corresponding soft Bellman equation
\citep{rust1987optimal,hotz1993conditional,aguirregabiria2010dynamic}. The same
softmax structure arises in modern maximum-entropy IRL by adding an entropy term
to the agent's objective, so actions with higher long-run utility are chosen
with higher probability rather than deterministically
\citep{ziebart2008maximum,ziebart2010causal,haarnoja2017reinforcement}. Even
under this structure, rewards remain identified only up to an equivalence
class. Point identification therefore requires a normalization
\citep{rust1987optimal,geng2020deep}, such as setting the reward of a reference
action to zero.



Uncertainty quantification for reward-dependent summaries is central
to IRL applications. These summaries include normalized reward averages,
structural coefficients, and values of counterfactual decision rules
\citep{kalouptsidi2021identification,zielnicki2025value}. The problem is
challenging because rewards are latent, identified only after normalization, and
estimated through flexible nuisance components such as choice probabilities and
transition laws. Existing methods largely split in two: classical DDC methods
provide inference under restrictive parametric utility models
\citep[e.g.][]{rust1987optimal,hotz1993conditional,aguirregabiria2010dynamic},
whereas flexible machine-learning approaches to IRL generally focus on reward
estimation without valid uncertainty quantification for downstream functionals
\citep[e.g.][]{abbeel2004apprenticeship,ramachandran2007bayesian,ziebart2008maximum,
levine2011nonlinear,ho2016gail,fu2018airl,wulfmeier2015deep,geng2020deep}.
We close this gap by developing a semiparametric framework for
\emph{debiased inverse reinforcement learning}.

 

\subsection{Our contributions}
We study reward identification and inference in maximum-entropy IRL and
Gumbel-shock dynamic discrete choice models. Our main contributions are as follows:
\begin{enumerate}
\item We show that, in softmax IRL, the reward is recoverable from the
log-behavior policy under an appropriate normalization. More generally, a broad
class of reward-dependent estimands can be expressed as smooth functionals of
the log-behavior policy and transition dynamics, reducing reward-based
inference to inference on observable choice probabilities and the transition
kernel.

\item We establish pathwise differentiability and derive efficient influence
functions, thereby characterizing the semiparametric efficiency bounds for
these targets. Under general affine reward normalizations, we show that the
efficient influence function can be obtained directly from the unnormalized
case, yielding a general template for inference under normalization.

\item We construct automatic debiased machine-learning estimators
\citep{chernozhukov2022automatic, van2025automatic} that allow flexible
nuisance estimation while achieving \(\sqrt{n}\)-consistency, asymptotic
normality, and semiparametric efficiency.

\item We specialize the general theory to derive efficiency bounds and
efficient estimators for several key IRL/DDC targets, including fixed-policy
values, counterfactual softmax policy values, normalized-reward policy values,
and structural coefficients of a normalized reward.
\end{enumerate}

The paper proceeds as follows. Section~\ref{sec:ident} shows how policy-value
differences, and under normalization the reward itself, can be identified from
observed behavior. Sections~\ref{sec::EIFsub} and~\ref{sec::est} develop the
semiparametric theory and automatic debiased machine-learning estimators for a
broad class of smooth target functionals. Section~\ref{sec::applications}
specializes these results to the main policy-value targets.

\subsection{Related work}

\noindent \textbf{Dynamic discrete choice models.}
Classical DDC work studies identification and efficient estimation of
finite-dimensional utility parameters using nested fixed-point, conditional
choice probability, and simulation-based methods
\citep{rust1987optimal,hotz1993conditional,hotz1994simulation,aguirregabiria2010dynamic}.
In contrast, we fix the softmax/Gumbel-shock structure of
\citet{rust1987optimal} and target reward-dependent functionals under
nonparametric reward representations. To our knowledge, debiased estimation,
efficient influence functions, and efficiency bounds for such functionals have
not been developed in this literature. For identification, the closest
connection is the conditional choice probability approach of
\citet{hotz1993conditional}: under Gumbel shocks, log-odds of observed choices
identify differences in choice-specific value functions. Rather than using
log-odds for value-difference identification, we use the log behavior policy
together with the transition law to identify normalized rewards and
reward-dependent functionals. An expanded related work is provided in Appendix~\ref{app::relatedwork}.

\noindent \textbf{Inverse reinforcement learning.}
Classical IRL recovers parametric or linear rewards under hard optimality
assumptions \citep[e.g.][]{ng2000algorithms,abbeel2004apprenticeship}, while
more recent work uses maximum-entropy or adversarial formulations with flexible
function approximators
\citep[e.g.][]{ziebart2008maximum,ziebart2010modeling,levine2011nonlinear,ho2016gail,fu2018airl,geng2020deep}.
Several papers formalize the resulting reward ambiguity, including
\citet{fu2018airl}, Theorem~1 of \citet{cao2021identifiability}, and Lemma~2 of
\citet{geng2020deep}. Deep PQR \citep{geng2020deep} identifies normalized
rewards under anchor-action constraints using a construction tailored to that
setting: the reward is represented through the log behavior policy, an
anchor-specific Bellman fixed point, and a second-stage conditional-mean
regression. In contrast, our identification results express normalized rewards
in terms of standard primitive objects---the log behavior policy, transition
law, and policy-specific \(Q\)-functions---and apply to a broader class of
affine normalization constraints, with anchor-action constraints as a special
case. These works focus on reward recovery; our goal is inference for
reward-dependent functionals.
functionals.

\noindent \textbf{Off-policy evaluation and double reinforcement learning.}
Debiased and doubly robust methods for off-policy evaluation estimate policy
values when the reward is known
\citep[e.g.][]{jiang2016doubly,thomas2016data,tang2019doubly,
liu2018breaking,van2018online,kallus2020double,uehara2020minimax,shi2021deeply,
kallus2022efficiently,mehrabi2024off}.
Double Reinforcement Learning extends this perspective to more general
functionals of the \(Q\)-function \citep{kallus2020double,kallus2022efficiently,van2025automaticDRL}.
Our framework is complementary: it treats the reward as unknown, recovers it
through the softmax behaviour model, and then performs semiparametric inference
for the induced reward-dependent targets.  

 

 

 
 
\section{Preliminaries}

\subsection{Problem setup} 
\label{sec::setup}

We consider an infinite-horizon, time-homogeneous Markov decision process (MDP) with state space \(\mathcal S\), action space \(\mathcal A\), and reference measure \(\mu\) on \(\mathcal S\). The initial state \(S_0\in\mathcal S\) is drawn from density \(\rho_0\) with respect to \(\mu\). At each time \(t\geq 0\), the agent occupies state \(S_t\), selects action \(A_t\in\mathcal A\) according to the behaviour policy \(\pi_0(\cdot\mid S_t)\), receives latent reward \(r_0^\dagger(A_t,S_t)\), and transitions to a next state \(S_{t+1}\in\mathcal S\) with conditional density \(k_0(\cdot\mid A_t,S_t)\) with respect to \(\mu\). In the IRL setting, the analyst observes states and actions, but not rewards. We write \(P_0\) for the law of the one-step observation \((S_0,A_0,S_1)\), and \(\mathbb P_0\) for the induced law of the full trajectory.



The inferential task is to learn reward-dependent summaries, such as policy
values, from observed state--action transitions. We observe
\(\mathcal{D}_n := \{(S_i, A_i, S_i')\}_{i=1}^n\), a collection of \(n\) i.i.d.\
samples from \(P_0\). Our analysis also accommodates temporal dependence through Markov-chain arguments under mixing conditions \citep{bibaut2021sequential,mehrabi2024off}.
  A fundamental challenge
is that \(r_0^\dagger\) is not point identified from state--action data alone:
many distinct rewards induce the same observable behaviour
\citep{ng1999policy,ng2000algorithms}. IRL methods address this by imposing an
optimality model linking the behaviour policy to the reward.

 
 

\subsection{Soft optimality and the soft Bellman equation}

We adopt the notion of \emph{soft optimality}, which arises in
Gumbel–shock-to-reward formulations \citep{rust1987optimal,hotz1993conditional,aguirregabiria2010dynamic}
and in maximum-entropy IRL \citep{ziebart2008maximum,ziebart2010modeling,haarnoja2017reinforcement}.
Soft optimality replaces deterministic reward maximization with an
entropy-regularized objective, yielding stochastic softmax decision rules. We
first recall the classical (hard)
value and $Q$-functions. For any policy $\pi:\mathcal{A}\times\mathcal{S}\to[0,1]$
and transition kernel $k:\mathcal{S}\times\mathcal{A}\times\mathcal{S}\to\mathbb{R}$,
let $\mathbb{E}_{k,\pi}$ denote expectation with respect to trajectories generated by
the MDP with transition kernel $k$, policy $\pi$, and initial state density $\rho_0$.

For any function $f:\mathcal{A}\times\mathcal{S}\to\mathbb{R}$, define the forward
operators
\[
\mathcal{P}_k f(a,s)
:= \int f(a,s')\,k(s'\mid a,s)\,\mu(ds'),
\qquad 
\pi f(s)
:= \sum_{a\in\mathcal{A}} \pi(a\mid s)\, f(a,s).
\]
For a discount factor $\gamma\in[0,1)$ and reward function
$r:\mathcal{A}\times\mathcal{S}\to\mathbb{R}$, the \emph{(hard) $Q$-function} is
\[
q_{r,k}^{\pi,\gamma}(a,s)
= \mathbb{E}_{k,\pi}\!\left[\sum_{t=0}^\infty \gamma^t r(A_t,S_t)
  \,\Big|\,A_0=a,\, S_0=s\right],
\]
the unique bounded solution to Bellman's fixed-point equation
\begin{equation}
q_{r,k}^{\pi,\gamma}(a,s)
= r(a,s) + \gamma\,\mathcal{P}_k V_{r,k}^{\pi,\gamma}(a,s),
\qquad 
V_{r,k}^{\pi,\gamma}(s) = \pi q_{r,k}^{\pi,\gamma}(s),
\label{eqn::hardQ}
\end{equation}
where $V_{r,k}^{\pi,\gamma}$ is the associated \emph{(hard) value function}
\citep{bellman1954theory}. Equivalently, it solves the inverse problem
\[
\mathcal{T}_{k,\pi,\gamma}(q_{r,k}^{\pi,\gamma}) = r,
\qquad  
\mathcal{T}_{k,\pi,\gamma}(f) := f - \gamma\,\mathcal{P}_k \pi f,
\]
where \(\mathcal{T}_{k,\pi,\gamma}\) is the Bellman operator.  Since \(\gamma < 1\), the operator \(\mathcal{T}_{k_0,\pi,\gamma}\) is invertible
on the Banach space of state–action functions equipped with the supremum norm
\citep{bellman1954theory}, and we denote its inverse by
\(\mathcal{T}_{k,\pi,\gamma}^{-1}\). Throughout, when the
discount factor is $\gamma$, we omit explicit dependence on $\gamma$ in our notation.  

Entropy regularization augments the reward $r_0(A_t,S_t)$ with an entropy bonus
$\tau\,\mathcal{H}(\pi(\cdot\mid S_t))$, where
$\mathcal{H}(\pi(\cdot\mid s)) := -\sum_{a} \pi(a\mid s)\log \pi(a\mid s)$ is the
Shannon entropy of the policy and $\tau \in (0,\infty)$ is a regularization
parameter. The value and $Q$-functions associated with this augmented reward are
referred to as the \emph{soft} value and $Q$-functions \citep{haarnoja2017reinforcement}.
Since rewards are identified only up to scale, we set $\tau=1$ without loss of
generality, interpreting rewards on a relative scale.

An agent is \emph{soft optimal} (with discount factor $\gamma$ and $\tau = 1$) if its policy
$\pi_0$ maximizes the \emph{soft value}, that is, the entropy-regularized
expected return,
\[
\pi \;\mapsto\;
\mathbb{E}_{k_0,\pi}\!\left[
  \sum_{t=0}^\infty \gamma^t 
  \{\, r_0^\dagger(A_t,S_t)
    + \mathcal{H}(\pi(\cdot\mid S_t)) \,\}
\right],
\]
over all policies~$\pi$ \citep{haarnoja2017reinforcement}.  The resulting soft-optimal policy is the softmax (multinomial logit)
distribution induced by the soft $Q$-function
\(r_0^\dagger(a,s) + \gamma\, v_0^\dagger(a,s)\):
\begin{equation}
\pi_0(a \mid s)
= \frac{\exp\!\left\{ r_0^\dagger(a,s) + \gamma\, v_0^\dagger(a,s) \right\}}
       {\sum_{a' \in \mathcal{A}}
        \exp\!\left\{ r_0^\dagger(a',s) + \gamma\, v_0^\dagger(a',s) \right\}},
\label{eqn::softmaxprob}
\end{equation}
where $v_0^\dagger$ satisfies the soft Bellman fixed-point system
\begin{align*}
v_0^\dagger(a,s)
&= \int V_0^\dagger(s')\, k_0(s' \mid a,s)\,\mu(ds'), \quad
V_0^\dagger(s)
= \log\!\left(
    \sum_{a' \in \mathcal{A}}
       \exp\!\bigl\{ r_0^\dagger(a',s) + \gamma\, v_0^\dagger(a',s) \bigr\}
\right),
\end{align*}
with $V_0^\dagger$ the soft value function of $\pi_0$ and $v_0^\dagger$ its
associated continuation value. Equivalently, in the maximum-entropy IRL formulation
\citep{ziebart2008maximum,ziebart2010modeling}, the agent’s policy maximizes
the entropy of the trajectory distribution subject to matching the observed
behaviour.\footnote{This yields a Boltzmann trajectory distribution,
\(\mathbb{P}_0^{\pi_0}((S_0,A_0,\ldots)=(s_0,a_0,\ldots))
\propto
\exp(\sum_{t=0}^\infty \gamma^t r_0^\dagger(a_t,s_t))\).}  
We take the discount factor \(\gamma\) as known. Identification of \(\gamma\)
from observed choices—requiring additional variation beyond that in our
setting---is analysed in \citet{abbring2020identifying}.

Soft optimality also arises in dynamic discrete-choice (DDC) models
\citep{rust1987optimal,hotz1993conditional,aguirregabiria2010dynamic}, where
the agent selects actions by maximizing a random utility $A_t \in \arg\max_{a \in \mathcal{A}}
\left\{ r_0^\dagger(a,S_t) + \gamma\, v_0^\dagger(a,S_t) + \varepsilon_t(a) \right\},$ where, in the DDC model, $v_0^\dagger(a,s)$ coincides with the classical
choice-specific continuation value and $\varepsilon_t(a)$ are i.i.d.\ Gumbel
(Type I extreme value) shocks \citep{rust1987optimal,pitombeira2024trajectory}.

 

\subsection{Notation}
\label{sec::notation}

We introduce the notation and regularity conditions used throughout.  Let
$\lambda := \# \otimes \mu$ denote the product measure on $\mathcal{A} \times
\mathcal{S}$, where $\#$ is the uniform measure on the finite action set
$\mathcal{A}$ and $\mu$ is a finite dominating measure on $\mathcal{S}$.  Thus,
for any measurable $B \subseteq \mathcal{A} \times \mathcal{S}$, $\lambda(B)
= \tfrac{1}{|\mathcal{A}|} \sum_{a \in \mathcal{A}}
  \int 1_B(a,s)\,\mu(ds).$
Let $\mathcal{H} := L^\infty(\lambda)$ denote the space of essentially bounded
reward functions, and let $\mathcal{K} \subset L^\infty(\mu \otimes \lambda)$
denote the space of essentially bounded transition kernels $k$ (conditional
densities of $S' \mid (A,S)$ with respect to~$\mu$) satisfying $c < k(s' \mid
a,s) < C$ for $\lambda$-a.e.\ $(a,s)$.  We equip $\mathcal{H}$ and $\mathcal{K}$
with the $L^2$ norms
\[
\|h\|_{\mathcal{H}}
:= \Bigg( \int \sum_{a \in \mathcal{A}} h(a,s)^2\,\mu(ds) \Bigg)^{1/2},
\qquad
\|k\|_{\mathcal{K}}
:= \Bigg( \int \sum_{a \in \mathcal{A}} k(s' \mid a,s)^2\,
            \mu(ds')\,\lambda(a,ds) \Bigg)^{1/2}.
\]
For readability, we keep only the most frequently recurring conventions here.
We write \(r_0^\dagger\) for the latent reward and
\(r_0(a,s):=\log \pi_0(a\mid s)\) for the pseudo-reward. We use \(q\) and
\(V\) for policy-specific \(Q\)- and state-value functions; when a scalar policy
value is needed, we write \(V_{r,k}(\pi;\gamma)\). 

We make the simplifying assumption that the initial state distribution,
behaviour policy, and transition dynamics are uniformly bounded away from zero
and infinity.
\begin{enumerate}[label=\textbf{(A\arabic*)}, ref=A\arabic*, series=posit]
\item \textit{(Bounded positivity).}\label{cond::boundedpositivity}
There exist constants \(c,C\in(0,\infty)\) and \(\eta\in(0,1)\) such that, for
\(\mu\otimes\lambda\)-a.e. $(s,a,s')\in\mathcal{S}\times\mathcal{A}\times\mathcal{S}$ ,$ c<\rho_0(s)<C,\qquad
\eta<\pi_0(a\mid s)$, $c<k_0(s'\mid a,s)<C.$
\end{enumerate}
All results are stated under Assumption~\ref{cond::boundedpositivity}. Let
\(\mathcal{M}\) denote the collection of distributions
\(P(ds,da,ds')=\rho_P(s)\pi_P(a\mid s)k_P(s'\mid a,s)\mu(ds')\lambda(da,ds)\)
satisfying this assumption. The boundedness requirement is imposed for
technical convenience and can in principle be relaxed using suitable moment
conditions on \(\rho_0\), \(\pi_0\), and \(k_0\).
 

 

 

 


%Under the above conditions, the measures $\lambda$ and $\lambda_0$ are equivalent with density $w := d\lambda_0 / d\lambda$ bounded as $0 < m \le w \le M < \infty$, so the corresponding $L^p$ norms and spaces are equivalent for all $p \in (0,\infty]$.

 

\section{Identification}
\label{sec:ident}

We first characterize reward indeterminacy under soft optimality, then show
that policy value differences remain identified, and finally impose a
normalization to identify the reward itself.

\subsection{Partial identification of the reward}
\label{sec::partialident}

Without additional constraints, the reward \(r_0^\dagger\) is only partially
identified: many rewards induce the same softmax policy for a soft-optimal
agent~\citep{cao2021identifiability}. Its scale is also unidentified; by
convention, we fix the entropy-regularization parameter at \(\tau=1\), which is
equivalent to fixing the softmax temperature
\citep{haarnoja2017reinforcement}. This subsection formalizes the resulting
partial identification. 

 
 
Formally, the partial identification of \(r_0^\dagger\) can be characterized
through the softmax likelihood \citep{rust1987optimal,ziebart2008maximum}.  The
observed policy \(\pi_0\) maximizes this soft Bellman--constrained likelihood,
so any reward \(r\) that induces the same softmax policy is also a solution of
the population M-estimation problem
\begin{equation}
\label{eqn::softbellmanloglik}
r \in \argmin_{r' \in \mathcal{H}} 
E_0\!\left[
  \log\!\Bigg(\sum_{a \in \mathcal{A}}
    \exp\{ r'(a,S) + \gamma v_{r',k_0}(a,S) \}\Bigg)
  - \{ r'(A,S) + \gamma v_{r',k_0}(A,S) \}
\right],
\end{equation}
where \(v_{r',k_0}\) satisfies the soft Bellman equation
\begin{equation}
v_{r',k_0}(a,s) 
= \int \log\!\left(\sum_{a' \in \mathcal{A}}
   \exp\!\big\{ r'(a',s') + \gamma v_{r',k_0}(a',s') \big\}\right)
\, k_0(s' \mid a,s)\,\mu(ds').
\label{eqn::softbellmanloss}
\end{equation}
Any reward function that induces the same behaviour as \(r_0^\dagger\) has the
potential-based reward--shaping form
\citep{ng1999policy,fu2018airl,geng2020deep}
\begin{equation}
r(a,s)
= r_0^\dagger(a,s)
  + c(s)
  - \gamma \int c(s')\, k_0(s' \mid a,s)\,\mu(ds'),
\label{eqn::equivclass}
\end{equation}
for a unique state-dependent potential \(c : \mathcal{S} \to \mathbb{R}\).  
Under this transformation, the soft \(Q\)-functions satisfy
\(q_{r,k_0}^{\mathrm{soft}}(a,s) 
 = q_{r_0^\dagger,k_0}^{\mathrm{soft}}(a,s) + c(s)\).

 

We begin with a trivial but informative solution to
\eqref{eqn::softbellmanloglik}: the log-behaviour policy
\(r_0(a,s) := \log \pi_0(a \mid s)\).  The soft Bellman constraint
\eqref{eqn::softbellmanloss} is satisfied with \(v_{r_0,k_0} \equiv 0\),
because the normalization \(\sum_{a' \in \mathcal{A}} \pi_0(a' \mid s) = 1\)
implies that the log-partition term vanishes,
\(\log \sum_{a' \in \mathcal{A}} \exp\{r_0(a',s)\} = 0\). 

\begin{lemma}[Trivial reward solution]
\label{theorem::trivialsol}
The log-behaviour policy \(r_0 = \log \pi_0\) is a solution to
\eqref{eqn::softbellmanloglik} and satisfies \eqref{eqn::equivclass} with
\(c(s) = -\,V_0^\dagger(s)\).
\end{lemma}
 

\subsection{Identification of policy values}
Although the reward is only partially identified, many quantities of practical
interest remain point identified.  Lemma~\ref{theorem::trivialsol} shows that
behaviour determines $r_0^\dagger$ only up to potential-based shaping.  Such
shifts change all policy values by the same constant and thus do not affect
policy comparisons.  Consequently, policy value differences remain identified,
and the log-behaviour-policy reward $r_0$ suffices for policy evaluation.  The
next result formalizes this fact.





\begin{theorem}[Behaviour policy identifies value differences]
\label{theorem::identvalue}
Suppose that \(r_0^\dagger\) solves \eqref{eqn::softbellmanloglik} for a
discount factor \(\gamma \in [0,1)\). Then, for any policy \(\pi\),
\[
q^{\pi,\gamma}_{r_0,k_0}(a,s)
= q^{\pi,\gamma}_{r_0^\dagger,k_0}(a,s) - V_0^\dagger(s),
\qquad
V_{r_0,k_0}(\pi;\gamma)
= V_{r_0^\dagger,k_0}(\pi;\gamma) - E_0[V_0^\dagger(S)].
\]
Hence, for any two policies \(\pi_1\) and \(\pi_2\), $V_{r_0,k_0}(\pi_1;\gamma) - V_{r_0,k_0}(\pi_2;\gamma)
=
V_{r_0^\dagger,k_0}(\pi_1;\gamma)
- 
V_{r_0^\dagger,k_0}(\pi_2;\gamma).$
\end{theorem}

 
\noindent A subtle but important point is that the identification of policy value
differences holds only for the discount factor $\gamma$ that appears in the
soft Bellman equation governing the agent’s behaviour.  To evaluate policies
under a different discount factor $\gamma'$, the reward must typically be
point identified.  Section~\ref{sec:normalization} shows how this can be
achieved via a normalization constraint.  Before turning to that result, we
illustrate Theorem~\ref{theorem::identvalue} in a counterfactual softmax
setting.



 
\begin{example}[Identification for counterfactual softmax policies] \label{ex:1}
Let $\mathcal{A}^\star \subseteq \mathcal{A}$ and $\tau^\star\in(0,\infty)$.
For $(r,k)\in\mathcal{H}\times\mathcal{K}$, define the counterfactual softmax policy
\begin{equation} \label{eqn::softmaxpolicycf}
\pi_{r,k}^\star(a\mid s)
\propto \mathbbm{1}\{a\in\mathcal{A}^\star\}
   \exp\!\bigl\{(r(a,s)+\gamma v_{r,k}^\star(a,s))/\tau^\star\bigr\},
\end{equation}
where \(v_{r,k}^\star\) solves the counterfactual soft Bellman equation
\begin{align}
\label{eqn::softbellmanCF}
v_{r,k}^\star(a,s)
\; &= \int \Phi_{r,k}^\star(s')\, k(s'\mid a,s)\,\mu(ds'),
\nonumber\\
\Phi_{r,k}^\star(s)
\; &:= \tau^\star \log\!\left(\sum_{a'\in\mathcal{A}^\star}
      \exp\!\left(\tfrac{1}{\tau^\star}(r(a',s)+\gamma v_{r,k}^\star(a',s))\right)\right).
\end{align}
Because $\pi_{r,k}^\star$ is invariant to potential-based reward shaping of the
form \eqref{eqn::equivclass}, and
Lemma~\ref{theorem::trivialsol} shows that $r_0$ and $r_0^\dagger$ differ by
such a transformation, we have
$\pi_{r_0^\dagger,k_0}^\star=\pi_{r_0,k_0}^\star$. Applying
Theorem~\ref{theorem::identvalue} to the pair of policies
$\pi_1=\pi_{r_0,k_0}^\star$ and $\pi_2=\pi_0$, the value difference
$V_{r_0^\dagger,k_0}(\pi_{r_0^\dagger,k_0}^\star;\gamma)
-
V_{r_0^\dagger,k_0}(\pi_0;\gamma)$
is identified using the log-behaviour-policy reward $r_0$ in place of
$r_0^\dagger$.
\qed
\end{example}




\subsection{Exact reward identification under normalization}
\label{sec:normalization}

The next theorem shows that, after normalization, the latent reward
\(r_0^\dagger\) is exactly recoverable from the log-behaviour policy reward
\(r_0\). We fix a reference policy \(\nu\) and a known state-only anchor
\(g:\mathcal{S}\to\mathbb{R}\), and impose that the average reward in state
\(s\) under \(\nu\) equals \(g(s)\):
\begin{equation}
\label{eqn::rdaggernorm}
    \nu r_0^\dagger(s)
    := \sum_{\tilde a \in \mathcal{A}} \nu(\tilde a \mid s)
       r_0^\dagger(\tilde a,s)
    = g(s),
    \quad P_0\text{-a.e.~} s.
\end{equation}


This constraint anchors the free state potential in \eqref{eqn::equivclass}
and covers common normalizations, including anchor-action, statewise
sum-to-zero, and value-based constraints
\citep{hotz1993conditional,bajari2010identification,
kallus2016revealed,geng2020deep}. For example, requiring
\(V_{r_0^\dagger,k}^{\nu,\gamma}(s)=h(s)\) for all \(s\) is equivalent, by the
Bellman equation, to choosing
\(g(s)=h(s)-\gamma\sum_{a\in\mathcal A}\nu(a\mid s)
\int h(\tilde s)\,k(d\tilde s\mid a,s)\).
The pair \((\nu,g)\) can also be chosen using auxiliary information, such as
outcomes that determine a reference reward level
\citep{kallus2018removing}. For instance, with auxiliary state--reward data
from the agent, one could take \(\nu=\pi\) and
\(g(s)=E_0\{r(A,S)\mid S=s\}\).


The following identification result shows that the normalized reward can be expressed
through the \(Q\)-function
\(q_{r_0-g,k_0}^{\nu,\gamma} = \mathcal{T}_{k_0,\nu,\gamma}^{-1}(r_0-g)\) of the
normalization policy \(\nu\).



\begin{theorem}[Normalization identifies the reward itself]
\label{theorem::identWithNormalization}
Let \(\nu\) be a policy.  
There exists a unique solution
\(r_{0,\nu,\gamma,g} \in \mathcal{H}\) 
to \eqref{eqn::softbellmanloglik}  that satisfies
\(\nu r_{0,\nu,\gamma,g}(s) = g(s)\) for \(P_0\text{-a.e.\ } s \in \mathcal{S}\).
Moreover,
\begin{align*}
    r_{0,\nu,\gamma,g}(a,s)
    &= g(s) + q_{r_0-g,k_0}^{\nu,\gamma}(a,s)
       - \sum_{\tilde a \in \mathcal{A}}
         \nu(\tilde a \mid s)\,
         q_{r_0-g,k_0}^{\nu,\gamma}(\tilde a,s), \\
    v_{0,\nu,\gamma,g}(a,s)
    &= \tfrac{1}{\gamma}\,\big(r_0(a,s) - g(s) - q_{r_0-g,k_0}^{\nu,\gamma}(a,s)\big),
\end{align*}
where  \(v_{0,\nu,\gamma,g}\) solves the soft Bellman equation with reward \(r_{0,\nu,\gamma,g}\). 
\end{theorem}

 


\noindent Thus, if \(r_0^\dagger\) satisfies \eqref{eqn::rdaggernorm}, it is
identified by \(r_{0,\nu,\gamma,g}\). Inference under normalization therefore
reduces to estimating the observable pseudo-reward \(r_0=\log\pi_0\) and the
Bellman quantity \(q_{r_0-g,k_0}^{\nu,\gamma}\), which motivates the next
section.





\section{Debiased Inverse Reinforcement Learning}

We now move from identification to inference. We first define a generic target
functional built from the pseudo-reward and transition kernel, then derive its
efficient influence function and bias expansion, and finally specialize these
objects to the motivating examples below.

\subsection{Statistical objective}
\label{sec::statobj}

The identification results in Section~\ref{sec:ident} imply that the
log-behaviour policy $r_0 = \log \pi_0$ can be used as a \emph{pseudo-reward}
for many IRL targets. Policy value differences are identified directly from
$r_0$ by Theorem~\ref{theorem::identvalue}, and
Theorem~\ref{theorem::identWithNormalization} shows that normalized reward
functionals can also be recovered from \((r_0,k_0)\). We therefore study
smooth real-valued functionals of the form
\[
\psi_0 := E_0[m(A,S,r_0,k_0)],
\]
where $(r,k)\mapsto m(\cdot,r,k)$ is a smooth map 
from $\mathcal{H}\times\mathcal{K}$ to $L^2(\lambda)$.
Formally, $\psi_0 = \Psi(P_0)$ corresponds to the parameter 
$\Psi:\mathcal{M}\to\mathbb{R}$ defined for $P\in\mathcal{M}$ by $\Psi(P) := E_P[m(A,S,r_P,k_P)],$ where $k_P(s'\mid a,s) := \tfrac{dP_{S'\mid A,S}}{d\mu}(s'\mid a,s)$ is the 
state transition kernel, and $r_P(a,s) = \log P(A=a \mid S = s)$.

\subsection{Motivating examples}
\label{subsec:examples}

Our framework accommodates a wide range of counterfactual targets in IRL.
We record the main targets used in the sequel: three policy-value parameters
and one structural normalized-reward parameter.  

\renewcommand{\theexample}{1a}
\begin{example}[Policy value]
\label{example::1a}
The policy value $\psi_0 := E_0[V_{r_0,k_0}^{\pi,\gamma}(S)]$ of a
policy $\pi$ with discount factor $\gamma$ corresponds to the map
$m(a,s,r,k) = V_{r,k}^{\pi,\gamma}(s)$.  
\qed
\end{example}

\noindent The same framework also covers interventions on the agent's
behaviour or environment.

\renewcommand{\theexample}{1b}
\begin{example}[Value of a soft-optimal agent]
\label{example::1b}
 Consider the counterfactual softmax policy $\pi_{r_0,k_0}^\star$ defined in
\eqref{eqn::softmaxpolicycf}, obtained by optimizing the pseudo-reward $r_0$
under kernel $k_0$ with action set $\mathcal{A}^\star$ and entropy parameter
$\tau^\star$.  
Unlike Example~\ref{example::1a}, this policy depends on both $r_0$ and $k_0$.
The policy value $\psi_0$ fits our framework through the map
$m(a,s,r,k) := V_{r,k}^\star(s)$, where \(V_{r,k}^\star\) denotes the value
function of \(\pi_{r,k}^\star\).  
\qed
\end{example}

% JASA trim: smooth the compressed transition from Example 1b to Example 1c.
\noindent A closely related counterfactual keeps the pseudo-reward target fixed
but replaces the transition kernel by a known \(k^\star\), as in simulated or redesigned environments
\citep[e.g.][]{abbeel2004apprenticeship,ziebart2008maximum,wulfmeier2015deep,fu2018airl}.

\renewcommand{\theexample}{1c}
\begin{example}[Policy values in counterfactual environments]
The value of the soft-optimal agent that maximizes the pseudo-reward $r_0$ in an
environment with known transition kernel $k^\star$ corresponds to
$m(a,s,r,k) := V_{r,k^\star}^\star(s)$, where \(V_{r,k^\star}^\star\) denotes
the value function of \(\pi_{r,k^\star}^\star\).  
\qed
\end{example}

Under a normalization constraint, functionals of the true reward
\(r_0^\dagger\) also fall within our framework. In these settings, the target
is identified by the \(\nu\)-normalized reward
\(\bar r_0 := r_{0,\nu,\gamma,g}\).

\renewcommand{\theexample}{2}
\begin{example}[Policy values under reward normalization]
\label{example::norm}
 Suppose we are interested in
\(\psi_0 = E_0[\tilde m(A,S,\bar r_0,k_0)]\),
a feature of the $\nu$-normalized reward \(\bar r_0\).  
By Theorem~\ref{theorem::identWithNormalization}, this can be expressed as
$\psi_0 = E_0[m(A,S,r_0,k_0)]$, where
\[
m(a,s,r,k)
:= \tilde m(a,s,r_{r,k,\nu,\gamma,g},k), \qquad
r_{r,k,\nu,\gamma,g}(a,s)
:= g(s) + q_{r-g,k}^{\nu,\gamma}(a,s) - V_{r-g,k}^{\nu,\gamma}(s),
 \]
and $r_{r,k,\nu,\gamma,g} = g + (I-\nu)\,\mathcal{T}_{k,\nu,\gamma}^{-1}(r-g)$ denotes
the $\nu$-normalized reward with anchor \(g\). This includes fixed-policy
values or softmax values built with a target discount \(\gamma'\) that may
differ from the normalization discount \(\gamma\).  
\qed
\end{example}

\renewcommand{\theexample}{3}
\begin{example}[Structural coefficients under reward normalization]
\label{example::normcoef}
Let \(x:\mathcal A\times\mathcal S\to\mathbb R^p\) be a known bounded feature
map and write \(\bar r_0:=r_{0,\nu,\gamma,g}\). If
\(G_0:=E_0\{x(A,S)x(A,S)^\top\}\) is nonsingular, the coefficient in the linear
projection of the normalized reward onto \(x\) is
\[
\theta_0
:=
G_0^{-1}E_0\{x(A,S)\bar r_0(A,S)\}.
\]
For example, in Rust's bus replacement model, a coordinate of \(x(a,s)\) may
indicate replacement, in which case the corresponding coordinate of
\(\theta_0\) is a replacement-cost coefficient. \qed
\end{example}

\subsection{Statistical efficiency and functional bias expansion}
\label{sec::EIFsub}

A key step in our methodology is to establish the pathwise differentiability of 
\(\Psi\) and derive its efficient influence function (EIF).  
In semiparametric statistics, the EIF determines the generalized 
Cramér--Rao lower bound for \(\Psi\) and forms the basis for constructing 
debiased, efficient estimators \citep{bickel1993efficient, laan2003unified, van2011targeted, chernozhukov2018double}.



The next conditions ensure that \(E_0[m(A,S,r,k)]\) is first-order
expandable in \(r\) and \(k\). The \emph{reward tangent space} \(T_{\mathcal{H}}\)
is defined as the completion of \(\mathcal{H}\) under the inner product $\langle h_1,h_2\rangle_{\mathcal{H},0}
:=E_0[h_1(A,S)h_2(A,S)].$ We define the \emph{kernel tangent space at \(k\)} by
\[
T_{\mathcal{K}}(k)
:= 
\Big\{
  \beta
  :
  \int \beta(s',a,s)^2\,k(s' \mid a,s)\,\mu(ds') < \infty
  \ \text{and}\
  \int \beta(s',a,s)\,k(s' \mid a,s)\,\mu(ds') = 0 
  \ \text{for $\lambda$-a.e. } (a,s)
\Big\},
\]
and set \(T_{\mathcal K} := T_{\mathcal K}(k_0)\), equipped with the inner product  
\(\langle \beta_1,\beta_2\rangle_{\mathcal{K},0}
:= E_0[\beta_1(S',A,S)\beta_2(S',A,S)]\).  
Intuitively, \(T_{\mathcal{K}}\) consists of allowable infinitesimal kernel scores that preserve the normalization of conditional densities under \(k(\cdot \mid a,s)\mu(ds')\).



 


\begin{enumerate}[label=\textbf{(C\arabic*)}, ref=C\arabic*, series=cond]
\item (\textit{Functional differentiability}) \label{cond::boundedfunc} 
There exist Gâteaux derivatives
$\partial_r m(\cdot, r_0, k_0) : T_{\mathcal{H}} \to L^2(\lambda)$ 
and $\partial_k m(\cdot, r_0, k_0) : T_{\mathcal{K}} \to L^2(\lambda)$
at $(r_0, k_0)$ such that, for all $\alpha \in T_{\mathcal{H}}$ and $\beta \in T_{\mathcal{K}}$,
\[
   \frac{d}{dt}\, E_0\!\left[m\big(A,S,\, r_0 + t \alpha,\,(1+t\beta)k_0\big)\right]\Big|_{t=0}
   = E_0\!\left[\partial_r m(A,S,r_0,k_0)(\alpha) 
   + \partial_k m(A,S,r_0,k_0)(\beta)\right].
\]

\item (\textit{Continuity of derivatives}) \label{cond::boundedfunc2}
The linear functionals $E_0[\partial_r m(A,S,r_0,k_0)(\cdot)]$ and $E_0[\partial_k m(A,S,r_0,k_0)(\cdot)]$
are continuous, in the sense that
\[
\sup_{\alpha \in T_{\mathcal{H}}}
\frac{E_0[\partial_r m(A,S,r_0,k_0)(\alpha)]}{\|\alpha\|_{\mathcal{H}, 0}}
< \infty,
\qquad
\sup_{\beta \in T_{\mathcal{K}}}
\frac{E_0[\partial_k m(A,S,r_0,k_0)(\beta)]}{\|\beta\|_{\mathcal{K}, 0}}
< \infty.
\]
\item (\textit{Lipschitz continuity}) \label{cond::lipschitz} The map $(r,k) \mapsto E_0[m(A,S,r,k)]$ is locally Lipschitz continuous at $(r_0,k_0)$ with respect to $(\mathcal{H}, \|\cdot\|_{\mathcal{H}, 0}) \times (\mathcal{K}, \|\cdot\|_{\mathcal{K}})$.
\end{enumerate}
Conditions \ref{cond::boundedfunc}-\ref{cond::lipschitz} hold if the map
\((r,k)\mapsto E_0[m(A,S,r,k)]\) is Fréchet differentiable at \((r_0,k_0)\)
with respect to the product norm
\(\|h\|_{\mathcal{H},0} + \|g\|_{\mathcal{K}}\); see
Lemma~\ref{lemma::frechetimplieslip} in Appendix~\ref{appendix::EIFsub}.
For additive Fréchet derivatives in the kernel argument, we write
\(\addTK\) for the ambient normed linear space of additive kernel perturbations,
equipped with \(\|\cdot\|_{\mathcal K}\), and denote the corresponding
additive derivative by \(\adk\). Thus \(\partial_k\) is reserved for score
derivatives on \(T_{\mathcal K}\), while \(\adk\) acts on additive directions
\(g \in \addTK\).



 

The EIF depends on two additional nuisance functions: the Riesz representers of
the partial functional derivatives of \((r,k)\mapsto E_0[m(A,S,r,k)]\). Let
\(\alpha_0\in T_{\mathcal H}\) and \(\beta_0\in T_{\mathcal K}\) denote these
representers with respect to \(\langle\cdot,\cdot\rangle_{\mathcal H,0}\) and
\(\langle\cdot,\cdot\rangle_{\mathcal K,0}\), so that
\begin{align}
E_0[\partial_r m(A,S,r_0,k_0)(\alpha)]
&= \langle \alpha_0, \alpha \rangle_{\mathcal{H}, 0}, 
\label{eqn::rieszreward}
\\[1ex] 
E_0[\partial_k m(A,S,r_0,k_0)(\beta)]
&= \langle \beta_0, \beta \rangle_{\mathcal{K}, 0}.
\label{eqn::rieszkernel}
\end{align}
 
 


 


\begin{theorem}[Pathwise differentiability]  \label{theorem::EIF}
    Under \ref{cond::boundedfunc}-\ref{cond::lipschitz}, the parameter
    \(\Psi: \mathcal{M} \to \mathbb{R}\) is pathwise differentiable at \(P_0\).
    Its efficient influence function is
    \[
    \chi_0(s,a,s')
    := \alpha_0(a,s)
    - \sum_{\tilde a \in \mathcal{A}} \pi_0(\tilde a \mid s)\alpha_0(\tilde a, s)
    + \beta_0(s,a,s')
    + m(a, s, r_0, k_0)
    - \Psi(P_0),
    \]
    where \(\alpha_0\) and \(\beta_0\) are the Riesz representers defined in
    \eqref{eqn::rieszreward}--\eqref{eqn::rieszkernel}.
\end{theorem}
\noindent The EIF \(\chi_0\) determines the generalized Cramér--Rao lower
bound for \(\Psi\). Its kernel component follows from existing EIF results for
offline reinforcement learning. Specifically, if the rewards
\(\{r_0(A_i,S_i)\}_{i=1}^n\) were observed, then
\(\beta_0(S_i,A_i,S_i')+m(A_i,S_i,r_0,k_0)-\Psi(P_0)\) would be the
nonparametric EIF for the known-reward target
\(P\mapsto E_P\{m(A,S,r_0,k_P)\}\). Thus, in each IRL application, it remains
only to derive the reward-side Riesz representer \(\alpha_0\) for the map
\(r\mapsto E_0\{\partial_r m(A,S,r,k_0)\}\). For example, EIFs for policy
values are given in
\citet{van2018online,kallus2020double,kallus2022efficiently}, and
\citet{van2025automaticDRL} derives EIFs for general continuous linear
functionals of the \(Q\)-function. Although those works allow additive reward
noise, the EIF for continuous linear functionals of the \(Q\)-function is
unchanged because the reward-noise score is orthogonal to the target-relevant
tangent space \citep{bickel1993efficient,laan2003unified}.


%The derivation of influence functions in IRL is greatly simplified by leveraging existing EIF results from offline RL. When the reward $r_0$ is known, i.e., when $\{r_0(A_i, S_i)\}_{i=1}^n$ are observed, the term $\beta_0(s,a,s') + m(a, s, r_0, k_0) - \Psi(P_0)$ coincides with the nonparametric EIF of the parameter $P \mapsto E_P[m(A, S, r_0, k_P)]$ based on the reward $r_0$. Thus, if the EIF of $P \mapsto E_P[m(A, S, r_0, k_P)]$ is available in the offline RL setting where rewards are observed, the component $\beta_0$ can be directly extracted. In particular, for the policy value of a known policy, the nonparametric EIF in the offline setting has been derived in \citet{van2018online, kallus2020double, kallus2022efficiently, van2025automaticDRL}, and more generally, the nonparametric EIF for continuous linear functionals of the $Q$-function---including the policy value as a special case---is provided in \citet{van2025automaticDRL}. These works technically assume that noisy rewards are observed, i.e., $\{R_i = r_0(A_i, S_i) + \varepsilon_i\}_{i=1}^n$ with $E_0[\varepsilon_i \mid A_i, S_i] = 0$ almost surely; however, for continuous linear functionals of the $Q$-function, the EIF under the noiseless reward model coincides with that in the noisy case when the noise is zero (i.e., $R = r_0(A,S)$). Consequently, for such parameters, if the nonparametric EIF of $P \mapsto E_P[m(A, S, r_0, k_P)]$ is already known from offline IRL, then only the Riesz representer $\alpha_0$ of $r \mapsto E_0[\partial_r m(A, S, r, k_0)]$ remains to be determined in order to obtain an explicit formula for the EIF in IRL.



 
  

 
The EIF characterizes the first-order behaviour of plug-in estimators of $\psi_0$.
To formalize this, we introduce the score
\[
\varphi_{r,k}(s',a,s;\alpha,\beta)
:= \alpha(a,s)
   - \sum_{\tilde a\in\mathcal{A}} \exp\{r(\tilde a,s)\}\,\alpha(\tilde a,s)
   + \beta(s',a,s),
\]
which provides the linear approximation of $E_0[m(A,S,r,k)]$ around $(r_0,k_0)$.
To analyze the bias of this approximation, define the remainder
\[
\mathrm{Rem}_0(r_0,k_0;\, r,k)
:= E_0[m(A,S,r,k)] - E_0[m(A,S,r_0,k_0)]
   - E_0[\varphi_{r_0,k_0}(\cdot;\alpha_0,\beta_0)].
\]
The next theorem establishes a von Mises expansion for this decomposition \citep{mises1947asymptotic}.

\begin{theorem}[von Mises expansion]
\label{theorem::vonmises}
Assume \ref{cond::boundedfunc}. Let $r \in \mathcal{H}$, $k \in \mathcal{K}$, 
$\alpha \in T_{\mathcal{H}}$, and $\beta \in T_{\mathcal{K}}(k)$, and suppose 
$\|r\|_{L^\infty(\lambda)} + \|r_0\|_{L^\infty(\lambda)} < M$ for some 
$M \in (0,\infty)$. Then
\begin{align*}
&E_0\!\big[m(A,S,r,k)\big] - E_0\!\big[m(A,S,r_0,k_0)\big] 
\;+\; \int \varphi_{r,k}(s',a,s;\alpha,\beta)\, P_0(ds,da,ds') \\
&\qquad
= -\,\langle r - r_0,\, \alpha - \alpha_0 \rangle_{\mathcal{H}, 0}
\;+\; O\!\big(\|r - r_0\|_{\mathcal{H}}^{2}\big) \\
&\qquad\quad
+\; O\!\Big(
      \|\beta_0 - \beta\|_{L^2(P_0)}
      \,\Big\|\tfrac{k - k_0}{k_0}\Big\|_{L^2(P_0)}
    \Big)
\;+\; \operatorname{Rem}_0(r_0,k_0;\,r,k),
\end{align*}
where the implicit constants depend only on $M$.
\end{theorem}
\noindent It is often possible to show that the linearization remainder is second order,
for example $\operatorname{Rem}_0(r_0,k_0;\,r,k)
\lesssim
\|k - k_0\|_{\mathcal{K}}^{2} + \|r - r_0\|_{\mathcal{H}}^{2}.$ For Example~\ref{example::1a}, this follows from Lemma~\ref{lemma::bellmanfrechet} in Appendix~\ref{appendix::applications}.

 

%This characterization shows that the plug-in estimator \(\tfrac{1}{n}\sum_{i=1}^n m(A_i,S_i,r_n,k_n)\) of \(\Psi(P_0)\) can be corrected by adding an estimate of its bias, namely \(\tfrac{1}{n}\sum_{i=1}^n \varphi_{r_n,k_n}(S_i',A_i,S_i;\alpha_n,\beta_n)\). Later, we will use this insight to construct debiased machine learning estimators for $\psi_0$ based on one-step estimation \citep{bickel1993efficient}.


\subsection{Functionals of \texorpdfstring{$\nu$}{nu}-normalized reward}
\label{sec::eifnorm}
Suppose we have already derived the EIF of a functional of the pseudo-reward $r_0$ using Theorem~\ref{theorem::EIF}, and we now seek the EIF of the corresponding functional of the $\nu$-normalized reward $r_{0,\nu,\gamma,g}$ from Theorem~\ref{theorem::identWithNormalization}.  
The next result shows that the EIF can be obtained directly from the unnormalized case via a simple transformation of the associated Riesz representers.

Consider the functional
$\psi_0 = E_0[\tilde{m}(A,S,r_{0,\nu,\gamma,g},k_0)]$ from Example~\ref{example::norm}.  
Let $\tilde{\alpha}_0$ and $\tilde{\beta}_{0}$ be the Riesz representers of the partial derivatives of the map $(r,k) \mapsto E_0[\tilde{m}(A,S,r,k)]$ at $(r_{0,\nu,\gamma,g}, k_0)$, defined by the property that for all $\alpha \in T_{\mathcal{H}}$ and $\beta \in T_{\mathcal{K}}$,
\begin{align*}
      E_0[\partial_r \tilde{m}(A,S,r_{0,\nu,\gamma,g},k_0)[\alpha]]
      &= \langle \tilde{\alpha}_0, \alpha \rangle_{\mathcal{H}, 0}, \\
      E_0[\partial_k \tilde{m}(A,S,r_{0,\nu,\gamma,g},k_0)[\beta]]
      &= \langle \tilde{\beta}_{0}, \beta \rangle_{\mathcal{K}, 0}.
\end{align*}
The Riesz representers are defined analogously to \eqref{eqn::rieszreward} and \eqref{eqn::rieszkernel}, the only difference being that the partial derivatives are now evaluated at $(r_{0,\nu,\gamma,g}, k_0)$ rather than $(r_0, k_0)$.  
The following theorem assumes Conditions~\ref{cond::boundedfunc::norm}--\ref{cond::lipschitz::norm}, which adapt Conditions~\ref{cond::boundedfunc}--\ref{cond::lipschitz} to the map $\tilde{m}$ evaluated at $(r_{0,\nu,\gamma,g}, k_0)$.  
Since these conditions are direct analogues, they are stated in Appendix~\ref{appendix::conditionsEIFnorm}.



\begin{theorem}[Pathwise differentiability for normalized rewards]
\label{theorem::EIFnorm}
Suppose \ref{cond::boundedfunc::norm}-\ref{cond::lipschitz::norm}, and define $m(a,s,r,k) := \tilde{m}(a,s,r_{r,k,\nu,\gamma,g},k).$ Then the parameter $\Psi: P \mapsto E_P[m(A,S,r_P,k_P)]$
is pathwise differentiable with EIF $\chi_0$, given in Theorem \ref{theorem::EIF}, where
 \begin{align*}
	      \alpha_0(a,s)
	      &:= \tilde{\alpha}_0(a,s)
	      - \frac{\nu(a \mid s)}{\pi_0(a \mid s)}
	         \sum_{\tilde a \in \mathcal{A}} \pi_0(\tilde a \mid s)\,
	         \tilde{\alpha}_0(\tilde a,s); \\[6pt]
	      \beta_0(s' \mid a,s) 
	      &:= \tilde{\beta}_{0}(s' \mid a,s)
	         + \alpha_0(a,s)\,
	           \Big\{r_0(a,s) - g(s) + \gamma V_{r_0-g, k_0}^{\nu, \gamma}(s') 
	           - q_{r_0-g, k_0}^{\nu, \gamma}(a,s)\Big\}.
	\end{align*}
\end{theorem}
\noindent Equivalently, Theorem~\ref{theorem::EIFnorm} gives a conditional
influence operator for the normalized reward itself.  For any downstream
functional whose derivative with respect to the normalized reward has Riesz
representer \(w\in T_{\mathcal H}\), define
\[
\mathcal A_\nu w(a,s)
:= w(a,s)
- \frac{\nu(a\mid s)}{\pi_0(a\mid s)}
   \sum_{\tilde a\in\mathcal A}\pi_0(\tilde a\mid s)w(\tilde a,s)
\]
and
\[
\Delta_\nu(s',a,s)
:= r_0(a,s)-g(s)
   +\gamma V_{r_0-g,k_0}^{\nu,\gamma}(s')
   -q_{r_0-g,k_0}^{\nu,\gamma}(a,s).
\]
Then the contribution to the EIF coming only from estimating
\(r_{0,\nu,\gamma,g}\) is
\[
\mathcal I_\nu(w)(s',a,s)
:= \mathcal A_\nu w(a,s)
   -\sum_{\tilde a\in\mathcal A}\pi_0(\tilde a\mid s)
      \mathcal A_\nu w(\tilde a,s)
   + \mathcal A_\nu w(a,s)\Delta_\nu(s',a,s).
\]
Thus, to compute the EIF for any smooth functional
\(E_0\{\tilde m(A,S,r_{0,\nu,\gamma,g},k_0)\}\), one may first compute the
ordinary Riesz representers \((w,\tilde\beta)\) of \(\tilde m\) as if the
normalized reward were observed, and then use
\[
\chi_0
= \mathcal I_\nu(w)
  + \tilde\beta_0(S',A,S)
  + \tilde m(A,S,r_{0,\nu,\gamma,g},k_0)-\psi_0 .
\]
This separates the reusable normalized-reward influence calculation from the
functional-specific derivative \(w\).

\noindent The simple structure of the EIF follows from the following representation.

\begin{lemma}[Representation of reward averages under normalization]
\label{lemma::linearfuncidentnorm}
For any \(w\in L^2(P_0)\),
\[
E_0\!\left[w(A,S)r_{0,\nu,\gamma,g}(A,S)\right]
=
E_0\!\left[
w(A,S)\left\{g(S)+r_0(A,S)
-\sum_{a\in\mathcal A}\nu(a\mid S)r_0(a,S)\right\}
\right].
\]
\end{lemma}

\noindent Hence, weighted averages of \(r_{0,\nu,\gamma,g}\), with possibly
signed weights \(w\), reduce to the reward-free term \(E_0\{w(A,S)g(S)\}\) and
an average of \(r_0\) after subtracting its \(\nu\)-mean at the same state.
Consequently, linear summaries such as the best linear predictor of
\(r_{0,\nu,\gamma,g}(A,S)\) can be identified from the pseudo-reward.

\section{Automatic debiased estimation and inference}
\label{sec::est}


\subsection{General form of the estimator}

Let $r_n$ and $k_n$ denote estimators of the pseudo-reward $r_0$
(the log-behaviour policy) and the transition kernel $k_0$, respectively.
The pseudo-reward $r_0$ can be learned from state–action data
$\{(S_i,A_i)\}_{i=1}^n$ via probabilistic classification, while $k_0$ may be
estimated using conditional density methods.  
A natural estimator of $\psi_0$ is the plug-in estimator
$n^{-1}\sum_{i=1}^n m(A_i,S_i,r_n,k_n)$.
Under flexible nuisance estimation, this plug-in approach can exhibit
substantial bias. We therefore use a one-step influence-function-based correction.

Let $\alpha_n \in \mathcal{H}$ and $\beta_n \in L^2(P_0)$ estimate the
representers $\alpha_0$ and $\beta_0$ from Section~\ref{sec::EIFsub}. We take
$\beta_n$ to be centered with respect to $k_n$, so that
$\int \beta_n(s',a,s)\,k_n(s' \mid a, s)\mu(ds')=0$; this can always be enforced post hoc. We propose the one-step estimator $\psi_n
:=
P_n\!\left\{
m(\cdot,r_n,k_n)
+
\varphi_{r_n,k_n}(\cdot\,;\alpha_n,\beta_n)
\right\},$ or, equivalently,
\begin{equation}
\psi_n
= \frac{1}{n}\sum_{i=1}^n m(A_i,S_i,r_n,k_n)
  + \frac{1}{n}\sum_{i=1}^n\!\left[
      \alpha_n(A_i,S_i)
      - \sum_{a\in\mathcal{A}}
          \exp\{r_n(a,S_i)\}\alpha_n(a,S_i)
      + \beta_n(S_i',A_i,S_i)
    \right].
\label{eqn::estgeneral}
\end{equation}
The second term applies the EIF from Theorem~\ref{theorem::EIF} to correct the
first-order bias of the plug-in estimator \citep{bickel1993efficient}. For policy
value parameters, the kernel representer \(\beta_0\) coincides with the usual
transition-kernel component of the EIF in off-policy evaluation, and can be
estimated using existing DRL methods. The reward representer can often be
derived analytically or estimated directly by \textit{Riesz regression}
\citep{chernozhukov2022automatic,van2025automatic}, using the population risk
\[
\alpha \mapsto E_0\left[\alpha(A,S)^2
-2\,\partial_r m(A,S,r_0,k_0)(\alpha)\right].
\]
Because \eqref{eqn::estgeneral} is agnostic to the choice of \(m\), we refer to
this procedure as an \emph{automatic} DML estimator.

When \(\psi_0\) depends on value- or \(Q\)-functions induced by \(k_0\), these
higher-level nuisances can often be estimated directly, avoiding explicit
estimation of \(k_0\). The next section uses this reparameterization for the main
policy-value examples.

 

 

\subsection{Asymptotic theory}

The next theorem establishes that the proposed estimator $\psi_n$ is asymptotically
linear with influence function $\chi_0$ and is therefore efficient for $\psi_0$.
We impose the following conditions.

\begin{enumerate}[label=\textbf{(C\arabic*)}, ref=C\arabic*, resume=cond]
\item (\textit{Sample splitting and boundedness}) \label{cond::samplesplit} \label{cond::bounded}
The estimators \(r_n\), \(k_n\), \(\alpha_n\), and \(\beta_n\) are obtained from data independent of \(\mathcal D_n\). Moreover, \(\bigl\|m(\cdot,r_n,k_n)+\varphi_{r_n,k_n}(\cdot\,;\alpha_n,\beta_n)\bigr\|_{L^2(P_0)}=O_P(1)\) and
\(\bigl\|m(\cdot,r_0,k_0)+\varphi_{r_0,k_0}(\cdot\,;\alpha_0,\beta_0)\bigr\|_{L^2(P_0)}=O(1)\).
 

\item (\textit{Nuisance consistency and estimation rates}) \label{cond::rates}
\begin{enumerate}[label=(\alph*)]
\item \label{cond::consistency} \(\bigl\|m(\cdot,r_n,k_n)+\varphi_{r_n,k_n}(\cdot\,;\alpha_n,\beta_n)-m(\cdot,r_0,k_0)-\varphi_{r_0,k_0}(\cdot\,;\alpha_0,\beta_0)\bigr\|_{L^2(P_0)}=o_P(1)\) 
    \item \label{cond::reward} \label{cond::representers} \(\|r_n-r_0\|_{\mathcal H}=o_P(n^{-1/4})\) and  \(\|r_n-r_0\|_{\mathcal H,r_0}\,\|\alpha_n-\alpha_0\|_{\mathcal H,r_0}=o_P(n^{-1/2})\).
    \item \label{cond::beta} \label{cond::remainder} \(\big\|\tfrac{k_n-k_0}{k_0}\big\|_{L^2(P_0)}\|\beta_n-\beta_0\|_{L^2(P_0)}=o_P(n^{-1/2})\) and  \(\operatorname{Rem}_0(r_0,k_0;\,r_n,k_n)=o_P(n^{-1/2})\).
\end{enumerate}
\end{enumerate}


\begin{theorem}[Asymptotic linearity and efficiency]
\label{thm:asymptotic_linearity}
Suppose Conditions~\ref{cond::samplesplit}--\ref{cond::rates} hold.  
Then the autoDML estimator $\psi_n$ defined in \eqref{eqn::estgeneral} satisfies
\[
\psi_n - \psi_0
= \frac{1}{n}\sum_{i=1}^n \chi_0(S_i,A_i,S_i') + o_P(n^{-1/2}).
\]
Consequently, $\sqrt{n}\,(\psi_n - \psi_0) \rightsquigarrow N(0,\sigma_0^2)$,  
where $\sigma_0^2 = \operatorname{Var}_{P_0}\!\big(\chi_0(S,A,S')\big)$.
Thus $\psi_n$ is nonparametric efficient for $\psi_0$.
\end{theorem}

\noindent The conditions above are standard in debiased machine learning.  
Condition~\ref{cond::samplesplit} enforces independence between nuisance estimation and evaluation; this can be relaxed to cross-fitting, where the sample is partitioned and nuisance functions are estimated on held-out folds \citep{van2011targeted, chernozhukov2018double}.  
Condition~\ref{cond::bounded} guarantees uniform $L^2$ boundedness of the true and estimated uncentered influence functions.  
Condition~\ref{cond::consistency} requires consistency of this estimated uncentered influence function, typically ensured when the nuisance estimators are consistent.  
Condition~\ref{cond::rates} specifies the convergence rates needed for the von Mises remainder to be second order.  
In particular, part~(c) yields a form of partial double robustness: slower convergence of $\alpha_n$ can be offset by faster convergence of $r_n$.  
In our examples, the linearization remainder often satisfies $\operatorname{Rem}_0(r_0,k_0;\, r,k) = O\big(\|r - r_0\|_{\mathcal H}^2 + \|k - k_0\|_{\mathcal K}^2\big)$.


 

\section{Applications}
\label{sec::applications}

% JASA move: make the main-text versus appendix split explicit at the start of Applications.
This section specializes the generic theory to the three policy-value targets
developed in the main text: the value of a fixed policy, the value of a fixed
policy under reward normalization, and the value of a counterfactual
soft-optimal policy. The corresponding normalized softmax and structural coefficient
examples are collected in Appendix~\ref{app::additionalapplications}.

\subsection{Efficient influence functions for example parameters}
\label{section::EIFexamples}

  We begin with Example~\ref{example::1a}, which concerns the value of a fixed
policy \(\pi\). Our analysis relies on an overlap condition for the stationary
state--action distribution, ensuring that the operator
\(\mathcal{T}_{k_0,\pi}\) has a bounded inverse on \(L^2(\lambda)\). For a
given policy \(\pi\), let \(p_{k}^{\pi,\gamma}\) denote a stationary
state-\(\mu\)-density for the MDP with policy \(\pi\) and kernel \(k\). That
is, for every \(f \in L^\infty(\mu)\), $\int f(s')\, p_{k}^{\pi,\gamma}(s')\, \mu(ds')
=
\int\!\!\int f(s')\, k_\pi'(s' \mid s)\, \mu(ds')\,
p_{k}^{\pi,\gamma}(s)\,\mu(ds),$ where
\(k_{\pi}'(s' \mid s):=\sum_{a \in \mathcal{A}} \pi(a \mid s)\, k(s' \mid a,s)\).





 
\begin{enumerate}[label=\textbf{(D\arabic*)}, ref=D\arabic*, series=condexample]
    %\item (\textit{Solution boundedness}) It holds that $\esssup_{a,s} \left|\frac{\pi_0(a \mid s)}{\exp \{r_0(a, s)\}}\right| < \infty$. \label{cond::boundedreward}
    \item (\textit{Stationary overlap for $\pi$}) There exist constants $b, B \in (0,\infty)$ such that  $b \;\le\; \frac{p_{k_0}^{\pi,\gamma}}{\rho_0}(s) \cdot \frac{\pi(a \mid s)}{\pi_0(a \mid s)} \;\le\; B $ for $\lambda$-a.e. $(a,s).$ \label{cond::stationary}
\end{enumerate}
 



In what follows, denote the (unnormalized) discounted state–action occupancy density ratio by
\begin{equation}
    d_0^{\pi,\gamma}(a,s)
    := \rho_0^{\pi,\gamma}(s)\,\frac{\pi(a \mid s)}{\pi_0(a \mid s)}, 
    \qquad 
    \rho_0^{\pi,\gamma}(s)
    := \frac{1}{\rho_0(s)}\sum_{t=0}^\infty \gamma^t\,
       \frac{d\mathbb{P}_{k_0,\pi}}{d\mu}(S_t = s),
    \label{eqn::occupancy}
\end{equation}
where $\rho_0^{\pi,\gamma}$ is the associated state occupancy ratio.
This state–action occupancy ratio characterizes the policy value through the identity
$E_0[V_{r_0,k_0}^{\pi,\gamma}(S)] = E_0[d_0^{\pi,\gamma}(A,S)\,r_0(A,S)]$.
In this sense, it plays the same role in offline policy evaluation as the inverse 
propensity score does in causal inference.
 


 


 
 

 

 
 
\begingroup
 
   
\begin{theorem}[EIF for known policy]
\label{theorem::EIFvalue}
 
   Let  $m(a,s,r,k) := V_{r,k}^{\pi,\gamma}(s)$. Under \ref{cond::stationary}, the parameter $\Psi: P \mapsto E_P[V_{r_P, k_P}^{\pi,\gamma}(S)]$ is pathwise differentiable at $P_0$ with efficient influence function
    \begin{align*}
    \chi_0(s,a,s') 
    &= \rho_0^{\pi,\gamma}(s)\!\left(\frac{\pi(a\mid s)}{\pi_0(a\mid s)}-1\right) \\[4pt]
     & \quad + d_0^{\pi,\gamma}(a,s)\Big\{r_0(a,s) + \gamma\,V_{r_0,k_0}^{\pi,\gamma}(s')
     - q_{r_0,k_0}^{\pi,\gamma}(a,s)\Big\} \\[4pt]
     & \quad + V_{r_0, k_0}^{\pi,\gamma}(s) - \Psi(P_0).
    \end{align*}
% Furthermore,  $\operatorname{Rem}_0(r_0,k_0; r,k) = O(\|k - k_0\|_{\mathcal{K}}^{2} + \|k - k_0\|_{\mathcal{K}}\|r - r_0\|_{\mathcal{H}}),$
%with an implicit constant depending only on $\gamma$, $\mathcal{K}$, $b$, $B$, and $\max\{\|r\|_{L^\infty(\lambda)}, \|r_0\|_{L^\infty(\lambda)}\}$.

\end{theorem}
\endgroup
 
\noindent In Theorem~\ref{theorem::EIFvalue}, the representers from 
Theorem~\ref{theorem::EIF} take the explicit forms
$\alpha_0(a,s) = \rho_0^{\pi,\gamma}(s)\,\tfrac{\pi(a\mid s)}{\pi_0(a\mid s)}$
and 
$\beta_0(s,a,s')
= d_0^{\pi,\gamma}(a,s)\{r_0(a,s) + \gamma\,V_{r_0,k_0}^{\pi,\gamma}(s') 
- q_{r_0,k_0}^{\pi,\gamma}(a,s)\}$.  
Thus the EIF combines a density-ratio term from the pseudo-reward with the
familiar temporal-difference term from off-policy evaluation
\citep{van2018online, kallus2020double, kallus2022efficiently, van2025automaticDRL}.  
%More generally, terms of the form ``weight $\times$ Bellman residual'' appear in EIFs of continuous linear functionals of the $Q$-function in offline RL \citep{van2025automaticDRL}.


 
We next consider Example~\ref{example::1b}, concerning the value of the
counterfactual softmax policy \(\pi_{r_0,k_0}^\star\). Here \(V^\star\)
denotes the policy value, \(v^\star\) the continuation value, and
\(\Phi^\star\) the log-sum-exp state value, where
\[
\Phi_{r_0,k_0}^\star(s)
:= \tau^\star \log\!\Biggl(\sum_{\tilde a \in \mathcal{A}^\star}
\exp\!\Bigl(\tau^{\star-1}\{r_0(\tilde a,s)
+ \gamma\,v_{r_0,k_0}^\star(\tilde a,s)\}\Bigr)\Biggr).
\]
We also define the conditional discounted state-occupancy ratio
\[
\rho_{r_0,k_0}^\star(s'\mid a,s)
:= \frac{1}{\rho_0(s')}\sum_{t=0}^\infty \gamma^t\,
\frac{\mathbb{P}_{k_0,\pi_{r_0,k_0}^\star}(S_t\in ds' \mid A_0=a,S_0=s)}{\mu(ds')},
\]
with marginal \(\rho_{r_0,k_0}^\star(s):=E_0[\rho_{r_0,k_0}^\star(s\mid A,S)]\)
and state-action ratio
\(d_{r_0,k_0}^\star(a,s):=\pi_{r_0,k_0}^\star(a\mid s)\pi_0(a\mid s)^{-1}
\rho_{r_0,k_0}^\star(s)\).
We also define the advantage-weighted state occupancy ratio
\begin{equation} \label{eqn::tilderho}
\widetilde{\rho}_{r_0,k_0}^\star(s)
:= E_0\!\big[d_{r_0,k_0}^\star(A,S)\{q_{r_0,k_0}^\star(A,S)
   - V_{r_0,k_0}^\star(S)\}\,\rho_{r_0,k_0}^\star(s\mid A,S)\big],
\end{equation}
with state--action version
$\widetilde{d}_{r_0,k_0}^\star(a,s)
:= \tfrac{\pi_{r_0,k_0}^\star(a\mid s)}{\pi_0(a\mid s)}\,
   \widetilde{\rho}_{r_0,k_0}^\star(s)$.
In the following condition, let
$p_{k_0, r_0}^{\star} := p_{k_0}^{\pi_{r_0, k_0}^\star, \gamma}$ denote the
stationary $\mu$-density under $\pi_{r_0, k_0}^\star$.



\begin{enumerate}[label=\textbf{(D\arabic*)}, ref=D\arabic*, resume=condexample]
    \item (\textit{Stationary overlap for $\pi_{r_0, k_0}^\star$}) There exist $b, B \in (0,\infty)$ such that $b \le \frac{p_{k_0, r_0}^{\star}(s)}{\rho_0(s)} \cdot \frac{\pi_{r_0, k_0}^\star(a \mid s)}{\pi_0(a \mid s)} \le B$ for $\lambda$-a.e.\ $(a,s)$. \label{cond::stationary3}
    \item (\textit{Lipschitz continuity of the soft Bellman fixed point}) There exists $L\in (0,\infty)$ such that, for all $k \in \mathcal{K}$, $\|v_{r_0,k}^\star - v_{r_0, k_0}^\star\|_{L^2(\lambda)} \le L \|k - k_0\|_{\mathcal{K}}$. \label{cond::lipschitzvalue}
\end{enumerate}

Condition~\ref{cond::lipschitzvalue} ensures pathwise differentiability of
$\Psi$ with respect to the kernel by controlling von Mises remainder terms of
the form $\|v_{r_0,k}^\star-v_{r_0,k_0}^\star\|_{L^2(\lambda)}$ in terms of
$\|k-k_0\|_{\mathcal K}$.
 
\begingroup
 
 
\begin{theorem}[EIF for softmax policy]
\label{theorem::EIFsoftmax}
Let $m(a,s,r,k) := V_{r,k}^\star(s)$. Under \ref{cond::stationary3} and \ref{cond::lipschitzvalue}, the parameter $\Psi: P \mapsto E_P[V_{r_P,k_P}^\star(S)]$ is pathwise differentiable at $P_0$ with efficient influence function
\begin{align*}
   \chi_0(s,a,s')
&= \left\{\frac{1}{\tau^\star}\,\widetilde{\rho}_{r_0,k_0}^\star(s)
    + \rho_{r_0,k_0}^{\star}(s)\right\}
   \left\{\frac{\pi_{r_0,k_0}^\star(a\mid s)}{\pi_0(a\mid s)} - 1\right\} \\[4pt]
& \quad + \frac{\gamma}{\tau^\star}\,\widetilde{d}_{r_0,k_0}^\star(a,s)
   \left\{\Phi_{r_0,k_0}^\star(s')-v_{r_0,k_0}^\star(a,s)\right\} \\[4pt]
& \quad + d_{r_0,k_0}^\star(a,s)\,\Big\{\,r_0(a,s) + \gamma\,V_{r_0,k_0}^{\star}(s') - q_{r_0,k_0}^\star(a,s)\,\Big\} \\[4pt]
& \quad + V_{r_0,k_0}^\star(s) - \Psi(P_0)
\end{align*}
%Furthermore, under only \ref{cond::stationary3},  $\operatorname{Rem}_0(r_0,k_0; r,k) 
%= O(\|k - k_0\|_{\mathcal{K}}^{2} 
%+ \|r - r_0\|_{\mathcal{H}}^{2} 
%+ \|v_{r_0,k}-v_{r_0,k_0}\|_{L^2(\lambda)}^{2}),$
%with an implicit constant depending only on $\gamma$, $\mathcal{K}$, $b$, $B$, %and $\max\{\|r\|_{L^\infty(\lambda)}, \|r_0\|_{L^\infty(\lambda)}\}$.

\end{theorem}
\endgroup
\noindent Relative to Theorem~\ref{theorem::EIFvalue}, this EIF adds the
advantage-weighted ratio $\widetilde{\rho}_{r_0,k_0}^\star$ and the Bellman
term
\(\tfrac{\gamma}{\tau^\star}\,\widetilde{d}_{r_0,k_0}^\star(a,s)
\{\Phi_{r_0,k_0}^\star(s') - v_{r_0,k_0}^\star(a,s)\}\).
Both terms arise because the target policy depends on the soft Bellman
solution \(v_{r_0,k_0}^\star\). 


We now turn to Example~\ref{example::norm}. Write
$d_{0,k_0}^{\pi,\gamma'}(a,s):=\rho_0^{\pi,\gamma'}(s)\,
\frac{\pi(a \mid s)}{\pi_0(a \mid s)}$ for the discounted state--action
occupancy ratio of \(\pi\) at discount \(\gamma'\), where
\(\rho_0^{\pi,\gamma'}\) is the associated state occupancy ratio. Define
$r_{0,\nu,\gamma,g}(a,s):=g(s)+q_{r_0-g,k_0}^{\nu,\gamma}(a,s)-
V_{r_0-g,k_0}^{\nu,\gamma}(s)$, and abbreviate
$q_{r_0,k_0,\nu}^{\pi,\gamma'}:=q_{r_{0,\nu,\gamma,g},k_0}^{\pi,\gamma'}$
and
$V_{r_0,k_0,\nu}^{\pi,\gamma'}(s):=\sum_{a\in\mathcal A}\pi(a\mid s)\,
q_{r_0,k_0,\nu}^{\pi,\gamma'}(a,s)$. In addition to the stationary overlap
condition from Condition~\ref{cond::stationary}, applied here with target
discount \(\gamma'\), we require stationary overlap for the normalization
policy \(\nu\):

\begin{enumerate}[label=\textbf{(D\arabic*)}, ref=D\arabic*, resume=condexample]
    \item (\textit{Stationary overlap for $\nu$}) There exist constants
    $b, B \in (0,\infty)$ such that
    $b \le \frac{p_{k_0}^{\nu,\gamma}(s)}{\rho_0(s)}
    \cdot \frac{\nu(a \mid s)}{\pi_0(a \mid s)} \le B$
    for $\lambda$-a.e.\ $(a,s)$. \label{cond::stationary2}
\end{enumerate}

\begingroup
\begin{theorem}[EIF for policy value under reward normalization]
\label{theorem::EIFvaluenorm}
Let $m(a,s,r,k)
:= \sum_{\tilde a \in \mathcal A}\pi(\tilde a \mid s)\,
q_{r,k,\nu}^{\pi,\gamma'}(\tilde a,s),$ where \(q_{r,k,\nu}^{\pi,\gamma'}\) is the \(Q\)-function of \(\pi\) at discount
\(\gamma'\) under the \(\nu\)-normalized reward \(r_{r,k,\nu,\gamma,g}\).
Under \ref{cond::stationary}, \ref{cond::stationary2}, and
\ref{cond::boundedfunc::norm}-\ref{cond::lipschitz::norm}, the parameter
\(\Psi: P \mapsto E_P[m(A,S,r_P,k_P)]\) is pathwise differentiable at \(P_0\)
with efficient influence function
\begin{align*}
\chi_0(s,a,s')
&= \rho_0^{\pi,\gamma'}(s)\left(\frac{\pi(a\mid s)}{\pi_0(a\mid s)}
   - \frac{\nu(a\mid s)}{\pi_0(a\mid s)}\right) \\[4pt]
&\quad + d_{0,k_0}^{\pi,\gamma'}(a,s)
   \left(1 - \frac{\nu(a \mid s)}{\pi(a \mid s)}\right)
   \bigl\{r_0(a,s) - g(s) + \gamma V_{r_0-g, k_0}^{\nu,\gamma}(s')
   - q_{r_0-g, k_0}^{\nu,\gamma}(a,s)\bigr\} \\[4pt]
&\quad + d_{0,k_0}^{\pi,\gamma'}(a,s)\,
   \bigl\{r_{0,\nu,\gamma,g}(a,s)
   + \gamma' V_{r_0, k_0,\nu}^{\pi,\gamma'}(s')
   - q_{r_0, k_0,\nu}^{\pi,\gamma'}(a,s)\bigr\} \\[4pt]
&\quad + V_{r_0,k_0,\nu}^{\pi,\gamma'}(s) - \Psi(P_0).
\end{align*}
\end{theorem}
\endgroup

\noindent Theorem~\ref{theorem::EIFvaluenorm} retains the same two-part
structure as Theorem~\ref{theorem::EIFvalue}, but adds a second Bellman term
for the normalization map. The first augmentation accounts for estimating the
normalized reward \(r_{0,\nu,\gamma,g}\) through the auxiliary value
\(q_{r_0-g,k_0}^{\nu,\gamma}\); the second is the usual off-policy correction
for evaluating \(\pi\) under that normalized reward.

 

Results for the additional examples from Section~\ref{subsec:examples},
namely softmax values under normalization and structural coefficients, are
collected in Appendix~\ref{app::additionalapplications}.

 



\subsection{Estimators and asymptotic theory}
\label{sec::examplesestimators}

We next give example-specific DML estimators that avoid direct estimation of
the transition kernel \(k_0\). Instead, they target higher-level nuisance
functions that yield doubly robust von Mises expansions and simpler rate
conditions under which Theorem~\ref{thm:asymptotic_linearity} applies. We focus
on Examples~ \ref{example::1a}, \ref{example::1b} and~\ref{example::norm}; additional estimators,
conditions, and asymptotic results are deferred to
Appendices~\ref{app:estimator:details} and~\ref{app::examplesestimators}.


We revisit Example~\ref{example::1a}. Let \(\pi_n:=\exp\{r_n\}\), let
\(q_n^{\pi,\gamma}\) estimate \(q_{r_0,k_0}^{\pi,\gamma}\), and define
\(V_n^{\pi,\gamma}(s):=\sum_{\tilde a\in\mathcal A}\pi(\tilde a\mid s)
q_n^{\pi,\gamma}(\tilde a,s)\). Let \(\rho_n^{\pi,\gamma}\) estimate the state
occupancy ratio \(\rho_0^{\pi,\gamma}\). Our DML estimator of
\(\psi_0=E_0[V_{r_0,k_0}^{\pi,\gamma}(S)]\) is
\begin{align*}
\psi_n
&:= \frac{1}{n}\sum_{i=1}^n V_n^{\pi,\gamma}(S_i)
  + \frac{1}{n}\sum_{i=1}^n
  \rho_n^{\pi,\gamma}(S_i)
  \frac{\pi(A_i\mid S_i)}{\pi_n(A_i\mid S_i)}
  \left\{
  r_n(A_i,S_i)+\gamma V_n^{\pi,\gamma}(S_i')
  -q_n^{\pi,\gamma}(A_i,S_i)
  \right\} \\
&\quad
  + \frac{1}{n}\sum_{i=1}^n
  \rho_n^{\pi,\gamma}(S_i)
  \left\{
  1-\frac{\pi(A_i\mid S_i)}{\pi_n(A_i\mid S_i)}
  \right\}.
\end{align*}
The nuisance \(q_{r_0,k_0}^{\pi,\gamma}\) can be estimated by fitted
\(Q\)-evaluation, i.e., iterative regression on Bellman targets
\citep{ernst2005tree,munos2005error,riedmiller2005neural,munos2008finite,mnih2013playing,van2025fitted}.
The ratio \(\rho_0^{\pi,\gamma}\) can be estimated by DICE-style,
moment-matching, or saddle-point methods
\citep{liu2018breaking,uehara2020minimax,nachum2019dualdice,zhang2020gendice,wang2023projected,van2025automaticDRL}.
For example, \(\rho_0^{\pi,\gamma}=\mathcal T_{k_0,\pi}(\widetilde\alpha_0)\)
\citep{van2025automaticDRL}, where
\[
\widetilde\alpha_0
:=
\argmin_{\alpha\in L^2(\mu)}
E_0\!\left[
\{\mathcal T_{k_0,\pi}(\alpha)(A_0,S_0)\}^2
-2\alpha(S_0)
\right].
\]
Thus one may estimate \(\rho_0^{\pi,\gamma}\) by replacing
\(\mathcal T_{k_0,\pi}\) with \(\mathcal T_{k_n,\pi}\), or by using an
equivalent min--max formulation
\citep{liu2018breaking,uehara2020minimax,dikkala2020minimax,kallus2020double}.
The next theorem gives rate conditions under which \(\psi_n\) is asymptotically
linear and efficient.

\begin{enumerate}[label=\textbf{(E\arabic*)}, ref=E\arabic*, series = efficiency1a]
  
    
    \item \label{assump:rates::policy} (\textit{Nuisance rates}) 
    $\bigl\{\|\rho_n^{\pi,\gamma} - \rho_0^{\pi,\gamma}\|_{L^2(\rho_0)} 
      + \|r_n - r_0\|_{\mathcal{H},r_0}\bigr\}\,
    \bigl\|
        \mathcal{T}_{k_0,\pi,\gamma}(q_n^{\pi,\gamma} - q_0^{\pi,\gamma})
    \bigr\|_{\mathcal{H},r_0}
    = o_p(n^{-1/2})$.
\end{enumerate}


 
\begin{theorem}[Asymptotic linearity for Example~\ref{example::1a}]
\label{theorem::ALex1}
Assume \ref{cond::stationary}, \ref{cond::rates}\ref{cond::reward}, and
\ref{assump:rates::policy}--\ref{assump:consistency::policy}.  
Then Theorem~\ref{thm:asymptotic_linearity} holds with influence function
$\chi_0$ given in Theorem~\ref{theorem::EIFvalue}.
\end{theorem}





For Example~\ref{example::1b}, let \(r_n\) estimate \(r_0\), let
\(v_n^\star\) estimate the soft Bellman fixed point
\(v_{r_0,k_0}^\star\), and define $\pi_n^\star(a\mid s)
\propto
1\{a\in\mathcal A^\star\}
\exp\!\left\{
\tfrac{1}{\tau^\star}\{r_n(a,s)+\gamma v_n^\star(a,s)\}
\right\}.$ Let \(q_n^\star\) estimate \(q_{r_0,k_0}^\star\), set $V_n^\star(s)
:=
\sum_{\tilde a\in\mathcal A^\star}
\pi_n^\star(\tilde a\mid s)q_n^\star(\tilde a,s),$
and let \(\widetilde\rho_n^\star\) and \(\rho_n^\star\) estimate
\(\widetilde\rho_{r_0,k_0}^\star\) and \(\rho_{r_0,k_0}^\star\). The resulting
estimator of \(\psi_0=E_0\{V_{r_0,k_0}^\star(S)\}\) is
\begin{align*}
\psi_n
&:=
\frac{1}{n}\sum_{i=1}^n V_n^\star(S_i)
+
\frac{1}{n}\sum_{i=1}^n
\rho_n^\star(S_i)
\frac{\pi_n^\star(A_i\mid S_i)}{\pi_n(A_i\mid S_i)}
\{r_n(A_i,S_i)+\gamma V_n^\star(S_i')-q_n^\star(A_i,S_i)\}
\\
&\quad+
\frac{\gamma}{n\tau^\star}\sum_{i=1}^n
\widetilde\rho_n^\star(S_i)
\frac{\pi_n^\star(A_i\mid S_i)}{\pi_n(A_i\mid S_i)}
\left[
\tau^\star
\log\!\left\{
\sum_{a\in\mathcal A^\star}
\exp\!\left(
\tfrac{1}{\tau^\star}\{r_n(a,S_i')+\gamma v_n^\star(a,S_i')\}
\right)
\right\}
-
v_n^\star(A_i,S_i)
\right]
\\
&\quad+
\frac{1}{n}\sum_{i=1}^n
\left\{
\tfrac{1}{\tau^\star}\widetilde\rho_n^\star(S_i)
+
\rho_n^\star(S_i)
\right\}
\left\{
\frac{\pi_n^\star(A_i\mid S_i)}{\pi_n(A_i\mid S_i)}-1
\right\}.
\end{align*}
The fixed point \(v_n^\star\) can be estimated by soft \(Q\)-iteration
\citep{haarnoja2017reinforcement,van2025stationary}, and \(q_n^\star\) by fitted
\(Q\)-evaluation under \(\pi_n^\star\). The ratio
\(\widetilde\rho_{r_0,k_0}^\star\) can be estimated by first estimating
\(\rho_{r_0,k_0}^\star(s'\mid a,s)\) and then applying
\eqref{eqn::tilderho}.  
\begin{enumerate}[label=\textbf{(E\arabic**)}, ref=E\arabic**, series=efficiency1b]
\item \label{assump:rates::softstar}(\textit{Nuisance rates})  
The following rates hold:
\begin{enumerate}[label=(\alph*)]
    \item \(\|v_n^\star - v_{r_0,k_0}^\star\|_{\mathcal{H},r_0}
    = o_p(n^{-1/4})\).
    
    \item \(\|\widetilde{\rho}_n^\star - \widetilde{\rho}_{r_0,k_0}^\star\|_{\mathcal{H},r_0}\,
    \bigl(\|r_n - r_0\|_{\mathcal{H},r_0} + \|v_n^\star - v_{r_0,k_0}^\star\|_{\mathcal{H},r_0}\bigr)
    = o_p(n^{-1/2})\).
    
    \item \(\bigl(\|\rho_n^\star - \rho_{r_0,k_0}^\star\|_{\mathcal{H},r_0}
      + \|r_n - r_0\|_{\mathcal{H},r_0}
      + \|v_n^\star - v_{r_0,k_0}^\star\|_{\mathcal{H},r_0}\bigr)\,
      \bigl\|\mathcal{T}_{r_0,k_0}^\star(q_n^\star - q_{r_0,k_0}^\star)\bigr\|_{\mathcal{H},r_0}
      = o_p(n^{-1/2})\).
\end{enumerate}
\end{enumerate}

 
 


\begin{theorem}[Asymptotic linearity for Example~\ref{example::1b}]
\label{theorem::ALex1b}
Assume \ref{cond::stationary3}, 
\ref{cond::rates}\ref{cond::reward}, and 
\ref{assump:rates::softstar}--\ref{assump:consistency::softstar}.  
Then Theorem~\ref{thm:asymptotic_linearity} holds with influence function 
$\chi_0$ given in Theorem~\ref{theorem::EIFsoftmax}.
\end{theorem}




 

 

 

For Example~\ref{example::norm}. Let \(\pi_n:=\exp\{r_n\}\), let
\(q_n^{\nu,\gamma}\) estimate \(q_{r_0-g,k_0}^{\nu,\gamma}\), and define
\(V_n^{\nu,\gamma}(s):=\sum_{\tilde a\in\mathcal A}\nu(\tilde a\mid s)\,
q_n^{\nu,\gamma}(\tilde a,s)\) and
\(r_{n,\nu,\gamma,g}(a,s):=g(s)+q_n^{\nu,\gamma}(a,s)-V_n^{\nu,\gamma}(s)\).
Let \(q_{n,\nu}^{\pi,\gamma'}\) estimate
\(q_{r_0,k_0,\nu}^{\pi,\gamma'}\), set
\(V_{n,\nu}^{\pi,\gamma'}(s):=\sum_{\tilde a\in\mathcal A}\pi(\tilde a\mid s)\,
q_{n,\nu}^{\pi,\gamma'}(\tilde a,s)\), and let \(\rho_n^{\pi,\gamma'}\)
estimate \(\rho_0^{\pi,\gamma'}\). The corresponding estimator of
\(\psi_0=E_0[V_{r_0,k_0,\nu}^{\pi,\gamma'}(S)]\) is
\begin{align*}
\psi_n
&:= \frac{1}{n}\sum_{i=1}^n V_{n,\nu}^{\pi,\gamma'}(S_i) \\
&\quad + \frac{1}{n}\sum_{i=1}^n \rho_n^{\pi,\gamma'}(S_i)\,
        \frac{\pi(A_i \mid S_i)}{\pi_n(A_i \mid S_i)}
        \bigl\{r_{n,\nu,\gamma,g}(A_i,S_i)
        + \gamma' V_{n,\nu}^{\pi,\gamma'}(S_i')
        - q_{n,\nu}^{\pi,\gamma'}(A_i,S_i)\bigr\} \\
&\quad + \frac{1}{n}\sum_{i=1}^n \rho_n^{\pi,\gamma'}(S_i)
        \left(
            \frac{\pi(A_i \mid S_i)}{\pi_n(A_i \mid S_i)}
            -
            \frac{\nu(A_i \mid S_i)}{\pi_n(A_i \mid S_i)}
        \right)
        \bigl\{r_n(A_i,S_i) - g(S_i)
        + \gamma V_n^{\nu,\gamma}(S_i')
        - q_n^{\nu,\gamma}(A_i,S_i)\bigr\} \\
&\quad + \frac{1}{n}\sum_{i=1}^n \rho_n^{\pi,\gamma'}(S_i)
        \left(
            \frac{\pi(A_i \mid S_i)}{\pi_n(A_i \mid S_i)}
            -
            \frac{\nu(A_i \mid S_i)}{\pi_n(A_i \mid S_i)}
        \right).
\end{align*}
 

 


 

\begin{enumerate}[label=\textbf{(E\arabic***)}, ref=E\arabic***, series=efficiency2]
    \item \label{assump:rates::policynorm} (\textit{Nuisance rates})
    The following rate conditions hold:
    \begin{enumerate}[label=(\alph*)]
        \item
        \(\bigl\{\|\rho_n^{\pi,\gamma'}-\rho_0^{\pi,\gamma'}\|_{L^2(\rho_0)}
        +\|r_n-r_0\|_{\mathcal{H},r_0}\bigr\}\,
        \bigl\|\mathcal{T}_{k_0,\pi,\gamma'}
        (q_{n,\nu}^{\pi,\gamma'}-q_{0,\nu}^{\pi,\gamma'})\bigr\|_{\mathcal{H},r_0}
        =o_p(n^{-1/2})\).
        \item
        \(\|r_n-r_0\|_{\mathcal{H},r_0}\,
        \Bigl\{
        \bigl\|\mathcal{T}_{k_0,\nu,\gamma}(q_n^{\nu,\gamma}-q_0^{\nu,\gamma})\bigr\|_{\mathcal{H},r_0}
        +\bigl\|\mathcal{T}_{k_0,\pi,\gamma'}
        (q_{n,\nu}^{\pi,\gamma'}-q_{0,\nu}^{\pi,\gamma'})\bigr\|_{\mathcal{H},r_0}
        \Bigr\}
        =o_p(n^{-1/2})\).
    \end{enumerate}
\end{enumerate}



 
\begin{theorem}[Asymptotic linearity for Example~\ref{example::norm}]
\label{theorem::ALex2}
Assume \ref{cond::stationary}, \ref{cond::stationary2},
\ref{cond::rates}\ref{cond::reward}, and
\ref{assump:rates::policynorm}--\ref{assump:consistency::policynorm}.
Then Theorem~\ref{thm:asymptotic_linearity} holds with influence function
\(\chi_0\) given in Theorem~\ref{theorem::EIFvaluenorm}.
\end{theorem}




 
\section{Simulation study}
\label{sec:simulation}

We evaluated the proposed estimators in a continuous-state treatment MDP with
two state variables, four actions, nonlinear rewards, and Gaussian state
transitions. The behaviour policy is the soft-optimal policy induced by the
data-generating reward and transition kernel, and all samples consist of
stationary one-step transitions. Throughout the simulation study, we treat the
transition kernel and discounted occupancy ratios as known, so that Bellman
equations and related population quantities can be computed exactly. The aim is
therefore not to study numerical error in solving Bellman equations, but to
isolate the statistical problem of reward estimation and the propagation of
reward-estimation error through policy and value-function operators. This
design allows us to study the three policy-value targets from
Section~\ref{sec::applications} within a common environment while retaining a
continuous state space and sufficient structure to compute oracle values on a
fine grid. All reported Monte Carlo summaries are based on $500$ replications
per design point.

Examples~\ref{example::1a} and~\ref{example::1b} illustrate complementary
finite-sample behaviour. Example~\ref{example::1a} is the most favourable
setting, and debiased IRL (D-IRL) exhibits near-nominal coverage together with a clear RMSE
improvement over the plug-in estimator at moderate and large sample sizes.
Example~\ref{example::1b} is more demanding because the target policy is
estimated from the recovered reward, so nuisance error propagates through the
soft Bellman map into both the target and the EIF correction. Even in that
setting, D-IRL reduces bias and RMSE at every sample size, although coverage
remains below nominal.
Table~\ref{tab:simulation-main} summarizes the results for these two examples.


\begin{table}[htb]
\centering
\caption{Examples~\ref{example::1a} and~\ref{example::1b}. Plug-in and D-IRL are
summarized by Monte Carlo bias, Monte Carlo standard deviation, and RMSE.
Coverage refers to the nominal $95\%$ Wald interval for D-IRL only.}
\label{tab:simulation-main}
\begingroup
\scriptsize
\setlength{\tabcolsep}{3pt}
\begin{tabular}{llrrrrrrr}
\toprule
Example & $n$ & Plug-in bias & Plug-in SD & Plug-in RMSE & D-IRL bias & D-IRL SD & D-IRL RMSE & Coverage \\
\midrule
1a & 1000 & -0.661 & 0.235 & 0.702 & 0.851 & 1.025 & 1.331 & 0.980 \\
1a & 2500 & -0.334 & 0.152 & 0.367 & 0.244 & 0.329 & 0.410 & 0.944 \\
1a & 5000 & -0.194 & 0.107 & 0.221 & 0.091 & 0.149 & 0.175 & 0.942 \\
1a & 10000 & -0.102 & 0.075 & 0.126 & 0.038 & 0.083 & 0.091 & 0.938 \\
1b & 1000 & 1.333 & 0.344 & 1.376 & 0.721 & 0.299 & 0.781 & 0.644 \\
1b & 2500 & 0.734 & 0.213 & 0.765 & 0.379 & 0.163 & 0.413 & 0.620 \\
1b & 5000 & 0.390 & 0.144 & 0.415 & 0.191 & 0.111 & 0.221 & 0.724 \\
1b & 10000 & 0.194 & 0.093 & 0.216 & 0.094 & 0.070 & 0.117 & 0.864 \\
\bottomrule
\end{tabular}
\endgroup
\end{table}

 
We report simulation results for Example~\ref{example::norm} in
Appendix~\ref{app::additionalapplications}. Inference is more challenging in
this setting because of the additional nuisance components, but coverage is
accurate when these components are well estimated.

\section{Conclusion}

We developed a semiparametric framework for inference in softmax IRL and dynamic
discrete choice models with nonparametric rewards. By expressing a broad class
of reward-dependent estimands as smooth functionals of the behaviour policy and
transition kernel, deriving their efficient influence functions, and
constructing automatic debiased ML estimators, the framework delivers valid and
computationally tractable inference under flexible reward representations. It
therefore connects classical inference for Gumbel-shock DDC models with modern
machine-learning-based IRL, and provides a reusable template for new
reward-dependent targets within the shared softmax structure.
Natural extensions include other shock structures, broader normalization
constraints, and nonhomogeneous or finite-horizon settings.




{
 \singlespacing 
\bibliographystyle{abbrvnat} % Or use 'abbrv' for another compact option
\bibliography{ref} % Assuming your .bib file is named 'ref.bib'
 }

\end{document}
